% from ega4-8-fr.tex
\oldpage[IV-3]{49}
\subsection{Traduction en termes de pro-objets}
\label{subsection:IV.8.13-fr}

La proposition suivante est essentiellement équivalente à (8.8.2, (i))~:

\begin{proposition}[8.13.1]
\label{IV.8.13.1-fr}
Soient $S$ un préschéma, $(X_\lambda,v_{\lambda\mu})$ un système projectif
filtrant de $S$-préschémas~; on suppose qu'il existe $\alpha$ tel que
$v_{\alpha\lambda}$ soit un morphisme affine pour tout $\lambda\geqslant
\alpha$ (ce qui entraîne (II, 1.6.2) que $v_{\lambda\mu}$ soit affine pour
$\alpha\leqslant\lambda\leqslant\mu$), de sorte que la limite projective
$X=\varprojlim_\lambda X_\lambda$ existe dans la catégorie des $S$-préschémas
(8.2.3). Soit $Y$ un $S$-préschéma, et, pour tout $\lambda\geqslant\alpha$,
soit
\[
e_\lambda:\operatorname{Hom}_S(X_\lambda,Y)
\longrightarrow\operatorname{Hom}_S(X,Y)
\]
l'application qui, à tout $S$-morphisme $f_\lambda:X_\lambda\to Y$, fait
correspondre $f=f_\lambda\circ v_\lambda$, où $v_\lambda:X\to X_\lambda$
est le morphisme canonique. La famille $(e_\lambda)$ est un système inductif
d'applications, qui définit donc une application canonique
\begin{equation}
\label{IV.8.13.1.1-fr}
\varinjlim_\lambda\operatorname{Hom}_S(X_\lambda,Y)
\longrightarrow\operatorname{Hom}_S(X,Y). \tag{8.13.1.1}
\end{equation}
Supposons $X_\alpha$ quasi-compact (resp. quasi-compact et quasi-séparé), et le
morphisme structural $Y\to S$ localement de type fini (resp. localement de
présentation finie). Alors l'application (8.13.1.1) est injective (resp.
bijective).
\end{proposition}

Posons en effet, pour $\lambda\geqslant\alpha$,
$Z_\lambda=Y\times_S X_\lambda$, de sorte que l'on a
$Z_\lambda=Z_\alpha\times_{X_\alpha}X_\lambda$. Posons de même
$Z=Y\times_S X=Z_\alpha\times_{X_\alpha}X$~; on sait alors (8.2.5) que, si
l'on pose aussi
\[
w_{\lambda\mu}=1\times v_{\lambda\mu}:Z_\mu\to Z_\lambda
\quad(\alpha\leqslant\lambda\leqslant\mu),
\qquad
w_\lambda=1\times v_\lambda:Z\to Z_\lambda
\quad(\alpha\leqslant\lambda),
\]
$Z$ est limite projective du système projectif $(Z_\lambda,w_{\lambda\mu})$
et les $w_\lambda$ sont les morphismes canoniques correspondants. Notons
d'autre part que le morphisme $Z_\alpha\to X_\alpha$ est localement de type
fini (resp. localement de présentation finie) (1.3.4 et 1.4.3). Enfin, on
sait que l'on a
\[
\operatorname{Hom}_S(X_\lambda,Y)
=\operatorname{Hom}_{X_\lambda}(X_\lambda,Z_\lambda)
\quad\hbox{et}\quad
\operatorname{Hom}_S(X,Y)=\operatorname{Hom}_X(X,Z)
\]
(I, 3.3.14). Il suffit maintenant d'appliquer (8.8.2, (i)) en y faisant
$X_\lambda=S_\lambda$ et en y remplaçant $Y_\lambda$ par $Z_\lambda$.

\begin{corollary}[8.13.2]
\label{IV.8.13.2-fr}
Avec les notations de (8.13.1), on suppose $X_\alpha$ quasi-compact et
quasi-séparé, et les $v_{\alpha\lambda}$ affines pour
$\alpha\leqslant\lambda$~; on suppose en outre que
$Y=\varprojlim_\rho Y_\rho$, où $(Y_\rho,t_{\rho\sigma})$ est un système
projectif filtrant de $S$-préschémas tel que, pour chaque $\rho$, le morphisme
structural $Y_\rho\to S$ soit localement de présentation finie. On a alors une
bijection canonique
\begin{equation}
\label{IV.8.13.2.1-fr}
\operatorname{Hom}_S(X,Y)
\xrightarrow{\sim}
\varprojlim_\rho\left(\varinjlim_\lambda
\operatorname{Hom}_S(X_\lambda,Y_\rho)\right). \tag{8.13.2.1}
\end{equation}
\end{corollary}

En effet, le fait que $Y$ soit limite projective des $Y_\rho$ entraîne en
particulier que l'application canonique
\[
\operatorname{Hom}_S(X,Y)\longrightarrow
\varprojlim_\rho\operatorname{Hom}_S(X,Y_\rho)
\]
est bijective~; et d'autre part, les hypothèses entraînent, pour chaque
$\rho$, l'existence d'une bijection canonique
\[
\operatorname{Hom}_S(X,Y_\rho)\xrightarrow{\sim}
\varinjlim_\lambda\operatorname{Hom}_S(S_\lambda,Y_\rho)
\]
en vertu de (8.13.1)~; d'où la conclusion.

\begin{env}[8.13.3]
\label{IV.8.13.3-fr}
Les résultats précédents permettent d'interpréter dans la théorie des
préschémas les notions de «~pro-variété~» ou de «~pro-schéma~» qui
s'introduisent dans certaines applications (par exemple dans la théorie du
corps de classes local suivant les idées de Serre [39] ou dans la théorie de
Néron sur la réduction des variétés
